
%      Unit 4 ,5
% 1.	Apply the concept of structured programming and modular design to develop a single library management system .
        justify how this approach improve readability and debugging. 

% 2.	Apply the step of software quality assurance and structured programming to create a basic student registration system .
%       Explain how your approach improve the quality and structure of the software. 

% 3. Demonstrate compliance width design and coding standard by creating a small program by creating a small program for student result
%   processing and documenting to naming convention. 

% 4. implement a real time health monitoring system using standard constructing practices and explain how the system minimize complexity and ensure maintainability . 
% 5. illustrate the application of software quality assurance  across sdlc phases of a banking application and provide suitable real time examples. 
% 6. solve the quality defect scenario in college erp system by applying sei cmm label 2 and label 3 practices explain the applied practices and improvement outcomes . 
% 7. explain the concept of top down programing to write a simple program for calculator explain how you broke down the problem into smaller parts . 6 marks
% 8. apply software quality assurance quality steps to a school result management system . explain how these steps help in improving the softwae quality. 
% 9. Demonstrate how top down programming help in building scalable software by developing a basic layout of an online shopping system. 
% 10. implement coding standards in a small c++ program that add and remove items from an inventory. explain how these standards enhance software maintainability. 
% 11. Illustrate the software quality assurance activities applied during the testing phase of a hospital management system project . 
% 12. Operate a sample software quality checklist to assist the quality of a student attendance portal. 
% 13. solve a real world quality problem using sqa principles and explain your approach and step taken .








----------------------------------------------------------------------------------------------------------------------------------------------

% Q.7.1Apply the Waterfall Model to a real-world software development scenario and explain its step-by-step implementation.
% OR
% Q.7.2Implement the Evolutionary Development Model in a scenario involving continuous customer feedback.

% Q.8.1Make use of a feasibility study to assess the viability of a new software application for a retail business. Explain how you would evaluate technical, operational, and economic feasibility.
% OR
% Q.8.2Use a DFD to represent the flow of information in a university management system. Demonstrate how DFDs can be used to capture both functional and non-functional requirements.

% Q.9.1Judge the appropriateness of using pseudo code over flowcharts in early design phases of a software project.
% OR
% Q.9.2Evaluate the advantages and limitations of function-oriented design versus object-oriented design in software development.

% Part C (Questions Only)
% Q.10.1Apply the Iterative Enhancement Model to develop an online e-commerce platform. Explain how the model would evolve through multiple iterations.
% OR
% Q.10.2Make use of the V-Model to manage the development of a mission-critical software for a financial institution. Discuss the role of testing in each phase of development.

% Q.11.1Make use of a feasibility study to evaluate the potential of a cloud-based inventory management system for a retail chain. Explain the steps involved in assessing technical, operational, and economic feasibility.
% OR
% Q.11.2Demonstrate a prioritized list of software requirements for an e-commerce website using the analysis phase of the requirement engineering process. Explain how to rank and justify the priority of each requirement to align with business objectives.

% Q.12.1Interpret a given flowchart and corresponding pseudo code to identify any design inefficiencies or logical inconsistencies.
% OR
% Q.12.2 Compare and evaluate different software measurement metrics with respect to their applicability in predicting software reliability and quality.





% 1) use an iterative enhancement model to develop an online e-commerce platform. explain how the model would evolve to multiple iterations.
% 2) implement a prototype based approach to develop a mobile application for restaurant reservation
% 3) Make a flow-chart and pseudocode to identify any design inefficiencies or logic inconsistencies
% 4) Justify the use of top-down design over bottom-up design for a library management software
% 5) make a DFD (data-flow diagram) for a hotel booking system and show how DFD can be used to model system processes and interaction, 
%      ensuring clarity and consistency in the system design. 
% Unit 3-
% 1) Evaluate the effectiveness of modularization in enhancing software reusability and maintainability in large scale software project
% 2) Critique the use of design-structured chart in representing the modular structure of a system. Explain their strength and limitations
% 3) Explain the role of function-oriented design versus object-oriented design in the development of real-time software system
% 4) Evaluate different software measurement matrix with respect to their application in predicting software reliability
-------------------------------------------------------------------------------------------------------------------------------------------------

